\documentclass[a4paper,11pt]{article}
\usepackage[utf8]{inputenc}
\usepackage[T1]{fontenc}
\usepackage[french]{babel}
\usepackage[left=2cm,right=2cm,top=2cm,bottom=2cm]{geometry}
\usepackage{xcolor}
\usepackage{hyperref}
\usepackage{fontawesome5}
\usepackage{parskip}

% --- Couleurs Cotral Lab ---
\definecolor{primary}{RGB}{0, 80, 140}
\hypersetup{colorlinks=true, urlcolor=primary}

\begin{document}

\noindent
\begin{minipage}[t]{0.5\textwidth}
    \textbf{NJIHA KUITO Brayanne Stella}\\
    Elève Ingénieure -- Master EEA\\[3pt]
    \faMapMarker\ Cergy (95), France\\
    \faPhone\ (+33) 6 52 31 08 46\\
    \faEnvelope\ \href{mailto:stellanjiha1@gmail.com}{stellanjiha1@gmail.com}
\end{minipage}
\hfill
\begin{minipage}[t]{0.4\textwidth}
    \raggedleft
    Cergy, le 1 juin 2026
\end{minipage}

\vspace{1.2cm}

\hfill
\begin{minipage}[t]{0.5\textwidth}
    \raggedleft
    \textbf{Cotral Lab / Comu Systems}\\
    Service Recrutement\\
    Emerainville (77), France
\end{minipage}

\vspace{1cm}

\noindent\textbf{Objet :} Candidature au poste d'Alternante Ingénieure R\&D -- Electronique Analogique \& Systèmes Embarqués -- Alternance 1 an

\vspace{0.8cm}

Madame, Monsieur,

La conception d'architectures électroniques analogiques, le prototypage de cartes et leur validation en laboratoire sont des activités que je pratique depuis ma première année à l'ENSEA. Contribuer aux projets R\&D de Comu Systems sur des systèmes de communication en milieu critique représente pour moi une opportunité unique d'appliquer ces compétences dans un contexte industriel exigeant, au service de professionnels dont la sécurité dépend de la fiabilité des équipements. C'est pourquoi je vous adresse ma candidature pour cette alternance à partir de septembre 2026, dans le cadre de mon Master 2 EEA à CY Paris Université.

Dans le cadre d'un projet de robot à navigation autonome à l'ENSEA, j'ai assuré de bout en bout la conception d'une carte électronique embarquée : saisie du schéma de principe sous KiCad, routage PCB en tenant compte des contraintes CEM et de puissance, puis assemblage, soudure de précision et mise au point du prototype en laboratoire. J'ai ensuite programmé les lois de commande moteur en C sur microcontrôleur STM32, avec debug et ajustement des paramètres jusqu'à la validation fonctionnelle complète. En parallèle, mes travaux de simulation sous LTSpice et PSpice sur des topologies de convertisseurs analogiques (Buck, Boost, onduleurs) m'ont développé une approche rigoureuse de la modélisation et de l'évaluation des performances des circuits.

Mon expérience en caractérisation CEM, où j'ai défini et mis en oeuvre des plans de mesures sur lignes microrubans et analysé les résultats pour identifier les axes d'amélioration, m'a appris à travailler avec méthode dans un contexte de validation exigeant. Je suis à l'aise avec les appareils de laboratoire (oscilloscope, GBF, analyseur de signaux) et la rédaction de rapports techniques.

Je suis disponible dès septembre 2026 pour une alternance d'un an et serais ravie de contribuer aux projets de conception et de validation de Comu Systems.

Je reste à votre disposition pour un entretien afin de discuter de ma candidature.

Je vous prie d'agréer, Madame, Monsieur, l'expression de mes salutations distinguées.

\vspace{0.8cm}

\textbf{NJIHA KUITO Brayanne Stella}\\
\textit{Elève Ingénieure ENSEA / CY Paris Université}

\end{document}
